% !TeX spellcheck = ru_RU
% !TEX root = vkr.tex

\section{Постановка задачи}
\label{sec:task}

Цель работы --- развить библиотеку \texttt{pysatl-mpest} для оценивания параметров и выбора структуры конечных смесей одномерных непрерывных распределений и подготовить её к выпуску.

На текущем обзорном этапе решаются следующие задачи:
\begin{enumerate}
    \item Сопоставить методы оценивания, инициализации и выбора структуры, а также программные инструменты по требованиям к однородным и разнородным смесям.
    \item Обосновать разделение ответственности библиотеки и вычислительного ядра, выбрать приоритетные средства развития и определить воспроизводимые сценарии проверки.
\end{enumerate}

Для дальнейшей разработки и экспериментальной проверки поставлены задачи:
\begin{enumerate}
    \item\label{task:reengineering}
          Завершить интеграцию оценщиков с \texttt{pysatl-core}, согласовав представление смеси, параметризации компонент и допустимые области параметров.
    \item\label{task:requirements}
          Доработать сменные стратегии оценивания и отдельные средства инициализации; обеспечить аналитические либо численные обновления выбранных подмножеств параметров, сохранение фиксированных параметров и диагностику недопустимых состояний.
    \item\label{task:structure-selection}
          Реализовать ограниченный поиск числа компонент и конфигураций допустимых семейств с оцениванием кандидатов и сравнением по информационным критериям; для выбранных вложенных моделей предусмотреть параметрический бутстреп отношения правдоподобий.
    \item\label{task:experiments}
          Воспроизвести выбранные фрагменты исследований: пример смеси Вейбулла из статьи о \texttt{mixR}, логнормально-паретовскую смесь с малым весом хвостовой компоненты и подгонку положительных смесей к датским страховым убыткам; проверить восстановление параметров и весов на синтетике и сопоставить аналитические и численные обновления по точности, времени и устойчивости.
    \item\label{task:release}\label{task:approbation}
          Подготовить выпуск библиотеки с документацией, примерами использования и сценариями воспроизведения, описав возможности и ограничения реализованных методов.
\end{enumerate}
